\documentclass{article}
\usepackage[T1]{fontenc}
\usepackage{lmodern}
\usepackage{graphicx}
\usepackage{amsmath,amssymb}
\usepackage[a4paper,margin=2cm]{geometry}
\usepackage[absolute,overlay]{textpos}
\setlength{\TPHorizModule}{1mm}
\setlength{\TPVertModule}{1mm}
\usepackage[indonesian]{babel}
\usepackage{hyperref}
\setlength{\emergencystretch}{2em}
% unit: Halaman 11

\begin{document}

% === Halaman 11 (python/native) ===

• 

Garis besar Algoritma Rijndael yang beroperasi pada blok  128-bit dengan kunci 128-bit 

adalah sebagai berikut (di luar proses pembangkitan round key): 

1. 

AddRoundKey: melakukan XOR antara state awal (plainteks) dengan cipher key. 

Tahap ini disebut juga initial round. 

2. 

Putaran sebanyak Nr – 1 kali. Proses yang dilakukan pada setiap putaran adalah: 

a. SubBytes: substitusi byte dengan menggunakan tabel substitusi (S-box). 

b. ShiftRows: pergeseran baris-baris array state secara wrapping. 

c. MixColumns: mengacak data di masing-masing kolom array state. 

d. AddRoundKey: melakukan XOR antara state sekarang round key. 

3. 

Final round: proses untuk putaran terakhir: 

a. SubBytes 

b. ShiftRows 

c. AddRoundKey

\end{document}
